\scanpage{34}
\unsure{entdt|entendu|ent.|ment}) les sommets s'identifient aux sections partielles sur des parties \inserted{de $J$} réduites à un élément -- la relation d'incidence est celle du prolongement des sections, ou encore l'inclusion des parties de $\Delta$.

Le groupe $G=\mathrm{Gl}(M)$ opère sur $\Delta$ par \unsure{autom|automs|automorphismes} d'octaèdre, on voit que $\mathfrak z=\{\pm1\}$ opère trivialement, \deleted{\unsure{d'où}} et $=$ que c'est là le noyau de l'opération, on a donc un hom injectif

\[
\underset{\substack{\rotatebox[origin=c]{90}{$\scriptstyle\simeq$}\\\mathrm{Gl}(M)/\{\pm1\}}}{\mathrm{Gl}(M)'}\longrightarrow\mathrm{Aut}_{\mathrm{oct}}(\Delta)
\]

\reversemarginpar\margin{\uline{NB} $\mathrm{Gl}(M)'\xrightarrow{\sim}\mathrm{GP}(M)$ (groupe projectif \unsure{en|de} \unsure{degré|dim|dimension|rang} $1$)}\normalmarginpar

qui est un isomorphisme, car les deux groupes sont de \deleted{\unsure{ce|un}} même ordre $24\times2=48$, donc on trouve

\[
\underset{\substack{\|\\\mathrm{Gl}(M)/(\mathfrak z=\{\pm1\})\\\rotatebox[origin=c]{90}{$\scriptstyle\simeq$}\\\deleted{\unsure{\mathrm P|\mathrm G}}\mathrm{GP}(M)}}{\mathrm{Gl}(M)'}
\xrightarrow{\sim}\mathrm{Aut}_{\mathrm{oct}}(\Delta)\simeq\mathrm{Aut}^+(\Delta)\times
\underbrace{\{1,\underset{\substack{\downarrow\\\text{antipodisme}}}{\underline a_\Delta}\}}_{\substack{\text{Centre}\\\text{de }\mathrm{Aut}_{\mathrm{oct}}(\Delta)}}
\tag*{(33)}
\]

Le centre de $\mathrm{Gl}(M)'$ \deleted{correspond} \inserted{n'est autre} d'ailleurs que le sous-groupe $\mathbb F_2^{\underline\omega}/\mathbb F_2\cdot\mathrm{id}_{M_0}$ de $\mathrm{Gl}(M)'$ (car le centre de $\unsure{\vert|\prime}\ \mathrm{Gl}(M_0)\simeq\mathfrak S_J$ est nul, donc le centre de $\mathrm{Gl}(M)'$ est contenu dans $\underline{\mathrm{End}}(M_0)/(\mathfrak z=\{\unsure{0_{M_0}|0_M|1_{M_0}},1_{M_0}\})\simeq M_0\times\mathbb F_2^{\underline\omega}/\mathbb F_2$, et s'identifie au sous-groupe des invariants sous $\mathfrak S_J\simeq\mathrm{Gl}_{\mathbb F_2}(M_0)$ de ce groupe, c'est $\mathbb F_2^{\underline\omega}/\mathbb F_2$.

La donnée d'un octaèdre $\Delta$ ayant $J$ comme ens des diagonales équivaut à celle d'un torseur
% archive sheet number: 33
